For any function $f:\RR^d\to\RR$ such that $Z_f\deq \int_{\RR^d} e^{-f(x)}dx<+\infty$, we define $\mu_f$ to be the distribution over $\RR^d$ with density $\mu_f(x)=\frac{1}{Z_f}e^{-f(x)}$.

For $\nu \ll \mu$, let $\Ren[\lambda]{\nu}{\mu} \deq \frac{1}{\lambda-1} \log \En_\mu[(\frac{d\nu}{d\mu})^\lambda]$ denote the R\'enyi divergence of order $\lambda > 1$. In addition, we define
\begin{align}
    \Dren{p}{q} \deq \E_p\prn[\big]{\frac{dp}{dq}}^\lambda - 1\,,\qquad
    \Dsys{p}{q}=\max\crl{\Dren{p}{q},\Dren{q}{p}}\,.
\end{align}


\subsection{Functional inequalities}


\begin{definition}[PI]\label{def:PI}
A probability distribution $\mu$ on $\R^d$ satisfies a Poincar\'e inequality (PI) with constant $C$ if for all smooth and compactly supported functions $\phi:\R^d\to \R$,
\begin{align*}
    \var_\mu(\phi)\leq C\En_\mu \nrm{\nabla \phi}^2\,.
\end{align*}
We let $\CPI(\mu)$ denote the smallest constant $C$ such that $\mu$ satisfies PI with constant $C$.
\end{definition}

\begin{definition}[LSI]\label{def:LSI}
A probability distribution $\mu$ on $\R^d$ satisfies a log-Sobolev inequality (LSI) with constant $C$ if for all smooth and compactly supported functions $\phi:\R^d\to \R$,
\begin{align*}
    \En_\mu\brk[\Big]{ \phi^2 \log \frac{\phi^2}{\En_\mu[\phi^2]}}\leq 2C\En_\mu \nrm{\nabla \phi}^2\,.
\end{align*}
We let $\CLSI(\mu)$ denote the smallest constant $C$ such that $\mu$ satisfies LSI with constant $C$.
\end{definition}

The following facts are standard, see~\citet{BGL14} or~\citet[Chapter 2]{Chewi26Book}.

\begin{lemma}[Bakry--\'Emery]\label{lem:LSI-SC}
If $\mu$ is $\alpha$-strongly log-concave, then $\CLSI(\mu)\leq \frac{1}{\alpha}$.
\end{lemma}

\begin{lemma}[Holley--Stroock perturbation principle]\label{lem:HS-perturb}
Let $\mu$ and $\nu$ be probability measures such that $\frac{1}{b}\leq \frac{d\nu}{d\mu}\leq b$ for some parameter $b>0$. Then $\CLSI(\nu)\leq b^2\,\CLSI(\mu)$.
\end{lemma}





\subsection{Technical lemmas}


\subsubsection{Tail estimates}

\begin{lemma}\label{lem:Gaussian-vMGF}
\begin{align*}
    \En_{W\sim \normal{0,\eta\Id}} \exp\prn*{\abs{\tri{W,v}}}\leq 2\exp\prn[\Big]{\frac{\eta\nrm{v}^2}{2}}\,.
\end{align*}
For $\lambda\leq \frac{1}{2\eta}$,
\begin{align*}
    \En_{W\sim \normal{0,\eta\Id}} \exp\prn[\Big]{\frac12\lambda\nrm{W+v}^2}\leq \exp\prn[\big]{d\lambda \eta+\lambda\nrm{v}^2}\,.
\end{align*}
Further, for $t\geq 0$,
\begin{align*}
    \PP_{x\sim \normal{v,\Sigma}}\prn*{\abs{x\tp Ax-v\tp A v-\tr(\Sigma A)}\geq t}&\leq 2\exp\prn*{ -\frac{1}{16\nrmop{\Sigma}}\min\crl*{\frac{t^2}{\nrmop{\Sigma}\nrmF{A}^2+\nrm{Av}^2}, \frac{t}{\nrmop{A}} }}\,.
\end{align*}
\end{lemma}
\begin{proof}
    \newcommand{\tSigma}{\wt{\Sigma}}
The first two statements are standard. For the third, for any $A\preceq \frac12\Sigma^{-1}$, we can calculate
\begin{align*}
    \En_{x\sim \normal{v,\Sigma}}\exp\prn[\Big]{\frac12\,x\tp Ax-\frac12\,v\tp A v}
    =\En_{W\sim \normal{0,\Sigma}} \exp\prn[\Big]{ v\tp A W+\frac12\, W\tp A W }
    =\sqrt{\frac{\det(\tSigma)}{\det(\Sigma)}}\exp\prn[\Big]{ \frac12\nrm{Av}_{\tSigma}^2 }\,,
\end{align*}
where $\tSigma^{-1}=\Sigma^{-1}-A$. Noting that $\Sigma^{1/2}\tSigma^{-1}\Sigma^{1/2}=\Id-\Sigma^{1/2}A\Sigma^{1/2}\succeq \frac12\Id$, we can bound
\begin{align*}
    \frac{\det(\tSigma)}{\det(\Sigma)}=\frac{1}{\det\prn*{\Id-\Sigma^{1/2}A\Sigma^{1/2}}}\leq \exp\prn[\big]{\tr(\Sigma^{1/2}A\Sigma^{1/2})+\nrmF{\Sigma^{1/2}A\Sigma^{1/2}}^2}\,,
\end{align*}
where we use $-\log(1-\lambda)\leq \lambda+\lambda^2$ for $\lambda\leq \frac12$. Therefore, we have shown
\begin{align*}
    \En_{x\sim \normal{v,\Sigma}}\exp\prn[\Big]{\frac12\,\prn*{x\tp Ax-v\tp A v-\tr(\Sigma A)}}
    \leq\exp\prn[\big]{\nrmF{\Sigma^{1/2}A\Sigma^{1/2}}^2+\nrm{Av}_{\Sigma}^2}\,, 
\end{align*}
as long as $A\preceq \frac12\Sigma^{-1}$. Then, by rescaling $A\leftarrow \lambda A$ and requiring $\abs{\lambda}\leq \frac{1}{2\nrmop{\Sigma^{1/2}A\Sigma^{1/2}}}$, we have
\begin{align*}
    \PP\prn*{\abs{x\tp Ax-v\tp A v-\tr(\Sigma A)}\geq t}\leq 2\exp\prn[\Big]{ \lambda^2\,\prn[\big]{\nrmF{\Sigma^{1/2}A\Sigma^{1/2}}^2+\nrm{Av}_{\Sigma}^2}-\frac12\lambda t }\,, \qquad \forall t>0\,.
\end{align*}
Suitably choosing $\lambda$ gives, for all $t\ge 0$,
\begin{align*}
    \PP\prn*{\abs{x\tp Ax-v\tp A v-\tr(\Sigma A)}\geq t}\leq&~ 2\exp\prn*{ -\min\crl*{\frac{t^2}{16(\nrmF{\Sigma^{1/2}A\Sigma^{1/2}}^2+\nrm{Av}_{\Sigma}^2)},\, \frac{t}{8\nrmop{\Sigma^{1/2}A\Sigma^{1/2}}} }} \\
    \leq&~ 2\exp\prn*{ -\frac{1}{16\nrmop{\Sigma}}\min\crl*{\frac{t^2}{\nrmop{\Sigma}\nrmF{A}^2+\nrm{Av}^2},\, \frac{t}{\nrmop{A}} }}\,.
\end{align*}
\end{proof}


\begin{lemma}\label{lem:Gaussian-s}
Suppose that $\eta>0$ and $0\leq \lambda\leq \frac{d^{1-s}}{4s\eta^s}$. Then, it holds that
\begin{align*}
    \En_{W\sim \normal{0,\eta\Id}} \exp\prn*{\lambda\nrm{W}^{2s}}\leq \exp\prn*{2(\eta d)^s \lambda}\,.
\end{align*}
\end{lemma}

\begin{proof}
We use the inequality $w^s\leq s\cdot\frac{w}{(\eta d)^{1-s}}+(1-s)\cdot (\eta d)^s$ for $w\geq 0$. Therefore, as long as $\frac{s\lambda}{(\eta d)^{1-s}}\leq \frac{1}{4\eta}$, it holds that
\begin{align*}
    \En\exp\prn*{\lambda\nrm{W}^{2s}}
    \leq \En\exp\prn[\Big]{\frac{s\lambda}{(\eta d)^{1-s}}\nrm{W}^{2}+(1-s)(\eta d)^s\lambda }
    \leq \exp\prn*{2s(\eta d)^s\lambda+(1-s)(\eta d)^s\lambda }
    \leq \exp\prn*{2(\eta d)^s \lambda}\,,
\end{align*}
where the second inequality follows from \cref{lem:Gaussian-vMGF}.
\end{proof}

\newcommand{\nrmexp}[1]{\nrm{#1}_{\psi_1}}
\begin{lemma}\label{lem:sub-exp-trunc}
For a random variable $Y$, we define the \emph{sub-exponential norm} of $Y$ as
\begin{align*}
    \nrmexp{Y}\deq \inf\crl{M>0: \En\exp(\abs{Y}/M)\leq 2}\,.
\end{align*} 
Then for any $s\geq 1$ and $t> 2s\nrmexp{Y}$, we have $\En(\abs{Y}^s-t^s)_+\leq 4t^se^{-t/\nrmexp{Y}}$. 
\end{lemma}
\begin{proof}
Fix any $M>\nrmexp{Y}$, and we only need to show $\En(\abs{Y}^s-t^s)_+\leq 4t^s e^{-t/M}$ for $t\geq 2sM$. Without loss of generality we assume $M=1$.  Note that we have $\PP(\abs{Y}\geq y)\leq 2e^{-y/M}$ for $y\geq 0$. Therefore, for $A\geq 2s$, 
\begin{align*}
    \En(\abs{Y}^s-A^s)_+
    \leq \En\abs{Y}^s \indic\crl{\abs{Y}\geq A}
    =s\int_A^{\infty} \PP(\abs{Y}\geq y)\,y^{s-1}\,dy
    \leq 2s\int_A^{\infty} e^{-y}\,y^{s-1}\,dy\leq 4sA^{s-1}e^{-A}\,.
\end{align*}
This is the desired upper bound. 
\end{proof}

\begin{lemma}\label{lem:Gaussian-Dcov}
Suppose $p=\normal{u,\eta\Id}$ and $q=\normal{v,\eta\Id}$. Then for $\delta\in(0,1)$, as long as $\eta\geq 8\log(1/\delta)\nrm{u-v}^2$, it holds that $p\prn[\big]{\frac{dp}{dq}\geq e}\leq \delta$.
\end{lemma}

\begin{lemma}\label{lem:Ito-Freedman}
    Suppose that $(B_t)_{t\ge 0}$ is the $d$-dimensional Brownian motion and $(J_t)_{t\ge 0}$ is a process of symmetric random matrices adapted to the filtration of $(B_t)_{t\ge 0}$.
Denote $Z_t\deq\int_0^t J_s\, dB_s$ and $M_t\deq \int_0^t J_s^2\,ds$. Then for any $R>0$, and distribution $u$ over $[0,T]$, 
\begin{align*}
    \bbP(\En_{t\sim u}\nrm{Z_t}^2\geq 4R\log(e/\delta))\leq 2\bbP(\tr(M_T)\geq R)+\delta\,.
\end{align*}
\end{lemma}

\begin{proof}
For any vector $v\in\RR^d$ and $\lambda\in\RR$, we know
\begin{align*}
    \En \exp\prn[\Big]{ \lambda \tri{v, Z_t}-\frac{\lambda^2}{2} \nrm{v}^2_{M_t} }\leq 1\,.
\end{align*}
Taking expectation over $v\sim \normal{0,\Id}$, it holds that
\begin{align*}
    \En \exp\prn[\Big]{ \frac{1}{2}\lambda^2 \nrm{Z_t}^2_{(\Id+\lambda^2 M_t)^{-1}}-\frac12\log\det(\Id+\lambda^2 M_t) }\leq 1\,.
\end{align*}
Note that $\nrm{Z_t}^2_{(\Id+\lambda^2 M_t)^{-1}}\geq \frac{1}{1+\lambda^2\nrmop{M_t}}\nrm{Z_t}^2$ and $\log\det(\Id+\lambda^2 M_t)\leq \lambda^2 \tr(M_t)$, we know
\begin{align*}
    \En \exp\prn[\Big]{ \frac{\lambda^2}{2(1+\lambda^2\nrmop{M_T})}\En_t\nrm{Z_t}^2-\frac{\lambda^2}2\tr(M_T) }
    \leq \En \En_{t}\exp\prn[\Big]{ \frac{\lambda^2}{2(1+\lambda^2\nrmop{M_t})}\nrm{Z_t}^2-\frac{\lambda^2}2\tr(M_t) }\leq 1\,.
\end{align*}
Therefore, by Markov's inequality, it holds that for any $\delta\in(0,1)$ and $\lambda>0$,
\begin{align*}
    \bbP\prn*{ \En_t\nrm{Z_t}^2\geq 2(1+\lambda^2\nrmop{M_T})\prn*{\lambda^{-2}\log(1/\delta)+\tr(M_T)} }\leq \delta\,.
\end{align*}
Choosing $\lambda^{-2}=R$ gives the desired upper bound.
\end{proof}



\subsubsection{Change of measure and tilts}

\begin{lemma}\label{lem:KL-triangle}

\begin{align*}
    \Dkl{\nu}{\muhat}\leq 2\Dkl{\nu}{\mu}+\log\prn*{1+\Dchis{\mu}{\muhat}}\,.
\end{align*}
\end{lemma}
\begin{proof}
We define $h=\log \frac{d\mu}{d\muhat}$. Then, we know that $\En_{\mu}[e^h]=1+\Dchis{\mu}{\muhat}$, and hence
\begin{align*}
    \Dkl{\nu}{\muhat}-\Dkl{\nu}{\mu}=&~\En_{\nu}\brk[\Big]{\log \frac{d\nu}{d\muhat}-\log \frac{d\nu}{d\mu}} 
    = \En_\nu[h] \\
    \leq&~ \Dkl{\nu}{\mu}+\log \En_{\mu}[e^{h}]
    =\Dkl{\nu}{\mu}+\log\prn*{1+\Dchis{\mu}{\muhat}}\,,
\end{align*}
where we use the Donsker--Varadhan variational inequality.
\end{proof}

\begin{lemma}\label{lem:Dchis-perturb}
Suppose that $\mu\propto pe^{h}$ and $\muhat\propto qe^{h}$, such that $\abs{h}\leq B$. Then it holds that
\begin{align*}
    \Dchis{\mu}{\muhat}\leq e^{4B}\Dchis{p}{q}\,.
\end{align*} 
\end{lemma}
\begin{proof}
By definition,
\begin{align*}
    1+\Dchis{\mu}{\muhat}
    =&~ \En_{x\sim \mu}\brk[\Big]{\frac{\mu(x)}{\muhat(x)}}
    =\frac{\En_q[e^h]}{(\En_p[e^h])^2}\cdot \En_{x\sim q}\brk[\Big]{ \frac{p(x)^2}{q(x)^2}\cdot e^{h(x)} }\,.
\end{align*}
Note that $\frac{p(x)^2}{q(x)^2}=\prn[\big]{\frac{p(x)}{q(x)}-1}^2+2\,\frac{p(x)}{q(x)}-1$, and hence
\begin{align*}
    \En_{x\sim q}\brk[\Big]{ \frac{p(x)^2}{q(x)^2}\cdot e^{h(x)} }
    =\En_{x\sim q}\brk[\Big]{ \prn[\Big]{\frac{p(x)}{q(x)}-1}^2 \cdot e^{h(x)} }+2\En_p[e^h]-\En_q[e^h]\,.
\end{align*}
Hence, we can rewrite
\begin{align*}
    \Dchis{\mu}{\muhat}
    =&~
    \frac{\En_q[e^h]}{(\En_p[e^h])^2}\cdot \En_{x\sim q}\brk[\Big]{ \prn[\Big]{\frac{p(x)}{q(x)}-1}^2 \cdot e^{h(x)} } - \prn[\Big]{\frac{\En_q[e^h]}{\En_p[e^h]}-1}^2 \\
    \leq&~ e^{4B}\En_{x\sim q}\brk[\Big]{ \prn[\Big]{\frac{p(x)}{q(x)}-1}^2  }
    =e^{4B}\Dchis{p}{q}\,.
\end{align*}
\end{proof}

\begin{lemma}\label{lem:Hels-var}
Suppose that $P,Q\in\DZ$. Then, for $h:\cZ\to[0,1]$,
\begin{align*}
    \En_P[h]\leq 3\En_Q[h]+4\Dhels{P}{Q}\,.
\end{align*}
Further, for any $f:\cZ\to [-1,1]$, it holds that
\begin{align*}
    \abs{\En_P[f]-\En_Q[f]}\leq 4\sqrt{\En_Q[f^2]\cdot \Dhels{P}{Q}}+4\Dhels{P}{Q}\,.
\end{align*}
\end{lemma}

\begin{proof}
We denote $P(\cdot)$ (resp.\ $Q(\cdot)$) to be the density function of $P$ (resp.\ $Q$). Then for any function $f:\cZ\to \RR$,
  \begin{align*}
    \abs{\En_P[f]-\En_Q[f]}^2=&~ \prn[\Big]{\int_{\cZ} f(z)\,(P(z)-Q(z))\, dz}^2 \\
    \leq&~\int_{\cZ} f(z)^2\,(\sqrt{P(z)}+\sqrt{Q(z)})^2 \, dz \cdot \int_{\cZ} (\sqrt{P(z)}-\sqrt{Q(z)})^2 \, dz \\
    \leq&~ 4\Dhels{P}{Q} \cdot \prn*{ \En_Q[f^2]+ \En_{P}[f^2]}\,.
  \end{align*}
  In particular, when $h:\cZ\to [0,1]$, the inequality above implies that 
  \begin{align*}
    \abs{\En_P[h]-\En_Q[h]}\leq 2\Dhel{P}{Q} \sqrt{ (\En_P[h]+\En_Q[h])}
    \leq \frac{1}{2}\,(\En_P[h]+\En_Q[h])+2\Dhels{P}{Q}\,,
  \end{align*}
  and hence it holds that $\En_P[h]\leq 3\En_Q[h]+4\Dhels{P}{Q}$. 

  Now, we can bound
  \begin{align*}
    \abs{\En_P[f]-\En_Q[f]}^2
    \leq&~ 4\Dhels{P}{Q} \cdot \prn*{ \En_Q[f^2]+ \En_{P}[f^2]} \\
    \leq&~ 16\Dhels{P}{Q} \cdot \prn*{ \En_Q[f^2]+ \Dhels{P}{Q}}\,.
  \end{align*}
   This gives the desired upper bound.
\end{proof}

\begin{lemma}\label{lem:KL-perturb}
Suppose that $\muhat\propto \mu e^{-h}$, where $\abs{h}\leq \B$. Then for any distribution $\nu$, it holds that
\begin{align*}
    \Dkl{\nu}{\muhat}\leq 4e^B\Dkl{\nu}{\mu}+4e^B\En_\mu[h^2]\,.
\end{align*} 
\end{lemma}
\begin{proof}
By definition,
\begin{align*}
    \Dkl{\nu}{\muhat}-\Dkl{\nu}{\mu}=&~\En_{\nu}[\log \mu(x)-\log \muhat(x)] \\
    =&~ \En_\nu[h]+\log \En_{\mu}[e^{-h}]
    \leq \En_\nu[h]-\En_\mu[h]+c_B\En_\mu[h^2]\,,
\end{align*}
where $c_B=\frac{e^B-B-1}{B^2}$. By \cref{lem:Hels-var}, it holds that
\begin{align*}
    \abs{\En_\nu[h]-\En_\mu[h]}\leq 4\sqrt{\En_\mu[h^2]\Dhels{\nu}{\mu}} + 4B\Dhels{\nu}{\mu} 
    \le 2\En_\nu[h^2] + (4B+2)\Dhels{\nu}{\mu}\,.
\end{align*}
Combining the inequalities above with $\Dhels{\nu}{\mu}\leq \frac12 \Dkl{\nu}{\mu}$ completes the proof.
\end{proof}






\begin{lemma}\label{lem:Renyi-to-diff}
For any $\ell>1$, it holds that
\begin{align}
    \Dsys[\ell]{\mu_f}{\mu_g}
    \leq \En_{x\sim \mu_f}\brk{e^{2\ell\,\abs{f(x)-g(x)}}-1}\,.
\end{align}
\end{lemma}
\begin{proof}
By definition, we can write
\begin{align*}
    \frac{\mu_f(x)}{\mu_g(x)}=e^{g(x)-f(x)}\En_{\mu_f}[e^{f-g}]\,.
\end{align*}
Therefore, we have
\begin{align*}
    1+\Dren[\ell]{\mu_f}{\mu_g}=&~\En_{\mu_f}\prn[\big]{\frac{\mu_f}{\mu_g}}^{\ell-1}
    = \prn*{\En_{\mu_f}[e^{f-g}]}^{\ell-1}\cdot \En_{\mu_f}[e^{(\ell-1)(g-f)}] \\
    \leq&~ \prn[\big]{\En_{\mu_f}e^{(\ell-1)\abs{f-g}}}^2
    \leq \En_{\mu_f}[e^{2\ell\abs{f-g}}]\,.
\end{align*}
Similarly,
\begin{align*}
    1+\Dren[\ell]{\mu_g}{\mu_f}=&~\En_{\mu_f}\prn[\big]{\frac{\mu_g}{\mu_f}}^{\ell}= \prn*{\En_{\mu_f}[e^{f-g}]}^{-\ell}\cdot \En_{\mu_f}[e^{\ell(f-g)}] 
    \leq \En_{\mu_f}[e^{-\ell(f-g)}]\cdot \En_{\mu_f}[e^{\ell(f-g)}]\\
    \leq&~ \prn[\big]{\En_{\mu_f}e^{\ell\abs{f-g}}}^2
    \leq \En_{\mu_f}[e^{2\ell\abs{f-g}}]\,.
\end{align*}
Combining both inequalities completes the proof.
\end{proof}















